\documentclass[11pt]{article}
\usepackage[utf8]{inputenc}
\usepackage{amsmath,amssymb}
\usepackage{geometry}
\geometry{margin=1in}
\setlength{\parindent}{0pt}
\setlength{\parskip}{6pt}

\begin{document}

Suppose we have $k$ coupons; in the graph, each node $u$ has adoption probability of $\alpha_u$, discarding probability $\beta_u$, and transfer probability $p_{u,v}$ for each out-neighbor $v$.

We first construct a possible world $W$, which realizes all randomness during the diffusion. A succinct representation of $W$ is that on each node $u$, we sample a random number
\[
r_u \sim \mathrm{Uniform}[0,1].
\]
If $r_u \le \alpha_u$, it means that node $u$ adopts the coupon when the coupon arrives at $u$. We also denote this event as $A_u$.  
If $\alpha_u < r_u \le \alpha_u + \beta_u$, it means that $u$ will discard the coupon. We denote this event as $D_u$.  
If $r_u > \alpha_u + \beta_u$, it means that $u$ will transfer the coupon to one of its out-neighbors. The actual neighbor is sampled according to parameter $p_{u,v}$. Let $v$ be the out-neighbor sampled, and $(u,v)$ be the sampled edge, which is commonly called a live edge.

Note that the above only considers one coupon. If we have $k$ coupons, they are independently transferred and adopted in the network. Thus, we need to use index $j$ to represent the randomness corresponding to the $j$-th coupon. That is, we will have
\[
r_u^{(j)}, \quad A_u^{(j)}, \quad D_u^{(j)}, \quad (u,v)^{(j)}.
\]

Therefore, one possible world is
\[
W = \bigl(W^{(1)}, W^{(2)}, \dots, W^{(k)}\bigr),
\]
or equivalently
\[
W^{(j)} = \{ r_u^{(j)}, A_u^{(j)}, D_u^{(j)}, (u,v)^{(j)} \}_{u \in V}.
\]
Each index $j$ corresponds the possible world for coupon $j$.

Suppose coupon $j$ is assigned to seed $s_j$. In the $j$-th coupon possible world, node $u$ adopts the $j$-th coupon if there is a path from $s_j$ to $u$ that only goes through live edges, and for every node
\[
w \text{ on the path } (w \neq u),
\]
node $w$ does not adopt the coupon nor discard the coupon, i.e.
\[
\neg A_w^{(j)} \wedge \neg D_w^{(j)},
\]
and for node $u$, it adopts the coupon, i.e. $A_u^{(j)}$. Denote this event as
\[
E_{s_j \rightarrow u}^{(j)}.
\]

But if there is a live edge from node $w$ pointing out to some other node, it automatically means that
\[
\neg A_w^{(j)} \wedge \neg D_w^{(j)},
\]
therefore, we do not need to explicitly say this condition. So the condition is only that there is a live-edge path from $s_j$ to $u$ and $A_u^{(j)}$.

The final adoption of $u$ is that $u$ adopts a coupon in at least one of the coupon possible worlds. That is
\[
A_u = \bigvee_{j=1}^k E_{s_j \rightarrow u}^{(j)}.
\]
When we want to clarify that this adoption is under the possible world $W$, we use $A_u(W)$. The $j$-th possible world is denoted as $W^{(j)}$. Note that here we use the simple model that the node's tranfer probabilities do not change after the node adopts the coupon. It will be more complicated and still open if the transfer probabilities change after the node adopts the coupon.

Let $\mathcal{R}(u, W^{(j)})$ denote the set of nodes $u$ can reach in the live-edge graph in world $W^{(j)}$. Let $RR(v, W^{(j)})$ be the reverse reachable set from root $v$ and the $j$-th world $W^{(j)}$, that is, the set of nodes that can reach $v$ in the live-edge graph of $W^{(j)}$.

Suppose the seed set is
\[
S = \{s_1, s_2, \dots, s_k\}.
\]

The adoption spread is
\[
\sigma(S) = \mathbb{E}_W \left[ \sum_{u} \mathbf{1}[A_u(W)] \right].
\]

If $k=1$, we only have one seed, then
\[
\sigma(\{s\})
= \sum_{v} \Pr\bigl( s \in RR(v, W) \wedge A_v \bigr)
= \sum_{v} \Pr(A_v)\Pr\bigl(s \in RR(v, W)\bigr).
\]

The second to last equality utilizes the fact that events $A_v$ and $s \in RR(v, W)$ are independent.

The above formula means that we could uniformly select a root node $v$, but the RR set from $v$ only has weight $\alpha_v$. This is equivalent to sampling root $v$ proportional to $\alpha_v$. Let $\mathcal{D}$ be the distribution of root node $v$ proportional to parameter $\alpha_v$. So the above would be the same as:
\[
\sigma(\{s\}) = \mathbb{E}_{v \sim \mathcal{D}} \bigl[ \mathbf{1}[s \in RR(v,W)] \bigr].
\]

The above means that when we have one seed node, we can randomly sample a root node proportional to node's adoption probability $\alpha_v$, then generate the reverse reachable set from $v$ based on the outgoing edge distribution. As we discussed, efficient sampling procedure exists.

Now, go back to $k$ seed nodes.

\[
\sigma(S)
= \sum_{v} \Pr\left(
\bigvee_{j=1}^k
\left(
s_j \in RR(v, W^{(j)}) \wedge A_v^{(j)}
\right)
\right).
\tag{*}
\]

The above formula means that, we can uniformly sample a root node $v$, then sample $k$ RR set from $v$, i.e. $RR(v,W^{(j)})$'s, and only keep those with event $A_v^{(j)}$. If none of the events $A_v^{(j)}$'s holds, then this sample an empty sample, resulting 0 contribution to the adoption spread.

What should we do if we want to avoid such zero contribution samples to make the sampling more efficient? It seems to require a sampling of events $A_v^{(j)}$'s for all $j$ conditioned on at least one of them is true. It may not be easy, unless we still independently sample $A_v^{(j)}$'s and just discard the case when all events are false. The probability that all $A_v^{(j)}$'s are false is
\[
\prod_{j=1}^k (1-\alpha_v),
\]
and hopefully it is small, so we do not waste too much time in doing such useless sampling.

To summarize, we have the following RR set sampling method:

1. Uniformly sample a node $v$.

2. For node $v$, sample $k$ random values in $[0,1]$ uniformly,
\[
r_v^{(1)}, \dots, r_v^{(k)},
\]
and determine the event $A_v^{(j)}$.

3. For $j$ from $1$ to $k$  
a. If $A_v^{(j)}$ is true, reverse generate RR set $RR(v,W^{(j)})$ from $v$ randomly. In this process, when we at a node $u$, we want to look at the in-neighbor $w$ of $u$; if $r_w^{(j)}$ has not be generated, then randomly generate $r_w^{(j)}$. Use $r_w^{(j)}$ to determine if there is an edge from $w$ to $u$ by checking if $r_w^{(j)}$ belongs to the interval dedicated to edge $(w,u)$. If so, add $w$ to the RR set; if not, ignore $w$ and go to the next in-neighbor of $u$. In traversal can be done in the DFS or BFS fashion.

How to use these RR sets is the next topic. We need to develop a greedy algorithm, similar to the IC model case. The difference is that the $j$-th seed node to be selected has their corresponding RR sets $RR(v,W^{(j)})$'s for all $v$.

Notes:

The formula (*) could be also used to prove the submodularity of $\sigma(S)$. One only needs to prove the submodularity of the function
\[
\mathbf{1}\!\left(
\bigvee_{j=1}^k
\left(
s_j \in RR(v,W^{(j)}) \wedge A_v^{(j)}
\right)
\right)
\]
for each $v$ and $W$.

RR set for coupon distribution  
Monday, March 25, 2024 4:33 PM

CouponDiffusion Page 1

1. The formula (*) could be also used to prove the submodularity of $\sigma(S)$.

2. The multiple RR set construction is similar to the RR set construction used in the following paper. The current one seems a bit more complicated due to the presence of events $A_v^{(j)}$'s.

\end{document}
